% Physical pages: 119--123 (printed pages 96--100).
% Figures: none.
% Uncertainties: none.

\chapterappendix{B}{Group Operators and Homotopy Classes}

In this appendix we shall prove the theorem stated in Section 2.7, that the
phases of the operators $U(T)$ for finite symmetry transformations $T$ may
be chosen so that these operators form a representation of the symmetry
group, rather than a projective representation, provided (a) the generators
of the group can be defined so that there are no central charges in the Lie
algebra, and (b) the group is simply connected. We shall also comment on
the projective representations encountered for groups that are not simply
connected, and their relation to the homotopy classes of the group.

To prove this theorem, let us recall the method by which we construct the
operators corresponding to symmetry transformations. As described in
Section 2.2, we introduce a set of real variables $\theta^a$ to parameterize
these transformations, in such a way that the transformations satisfy the
composition rule \eqref{eq:2.2.15}:
\begin{equation*}
  T(\bar\theta)T(\theta)=T\bigl(f(\bar\theta,\theta)\bigr).
\end{equation*}
We want to construct operators $U(T(\theta))\equiv U[\theta]$ that satisfy
the corresponding condition\footnote{Square brackets are used here to
distinguish $U$ operators constructed as functions of the group parameters
from those expressed as functions of the group transformations themselves.}
\begin{equation}
  U[\bar\theta]U[\theta]=U\bigl[f(\bar\theta,\theta)\bigr].
  \tag{2.B.1}\label{eq:2.B.1}
\end{equation}
To do this, we lay down arbitrary `standard' paths
$\Theta^a_{\theta}(s)$ in group parameter space, running from the origin to
each point $\theta$, with $\Theta^a_{\theta}(0)=0$ and
$\Theta^a_{\theta}(1)=\theta^a$, and define $U_{\theta}(s)$ along each such
path by the differential equation
\begin{equation}
  \frac{d}{ds}U_{\theta}(s)
    =it_aU_{\theta}(s)h^a{}_{b}\bigl(\Theta_{\theta}(s)\bigr)
      \frac{d\Theta^b_{\theta}(s)}{ds}
  \tag{2.B.2}\label{eq:2.B.2}
\end{equation}
with the initial condition
\begin{equation}
  U_{\theta}(0)=\mathbf{1},
  \tag{2.B.3}\label{eq:2.B.3}
\end{equation}
where
\begin{equation}
  [h^{-1}]^a{}_{b}(\theta)
    \equiv\left[
      \frac{\partial f^a(\bar\theta,\theta)}{\partial\bar\theta^b}
    \right]_{\bar\theta=0}.
  \tag{2.B.4}\label{eq:2.B.4}
\end{equation}
We are eventually going to identify the operators $U[\theta]$ with
$U_{\theta}(1)$, but first we must establish some of the properties of
$U_{\theta}(s)$.

In order to check the composition rule, consider two points $\theta_1$ and
$\theta_2$, and define a path $\gamma$ that runs from 0 to $\theta_1$
and thence to $f(\theta_2,\theta_1)$:
\begin{equation}
  \Theta^a_{\gamma}(s)\equiv
  \begin{cases}
    \Theta^a_{\theta_1}(2s),
      &0\leq s\leq \tfrac12,\\
    f^a\bigl(\Theta_{\theta_2}(2s-1),\theta_1\bigr),
      &\tfrac12\leq s\leq 1.
  \end{cases}
  \tag{2.B.5}\label{eq:2.B.5}
\end{equation}
At the end of the first segment, we are at
$U_{\gamma}(\tfrac12)=U_{\theta_1}(1)$. To evaluate
$U_{\gamma}(s)$ along the second segment, we need the derivative of
$f^a(\Theta_{\theta_2}(2s-1),\theta_1)$. For this purpose, we use the
fundamental associativity condition:
\begin{equation}
  f^a\bigl(f(\theta_3,\theta_2),\theta_1\bigr)
    =f^a\bigl(\theta_3,f(\theta_2,\theta_1)\bigr).
  \tag{2.B.6}\label{eq:2.B.6}
\end{equation}
Matching the coefficients of $\theta_3^c$ in the limit $\theta_3\to0$
yields the result:
\begin{equation}
  \frac{\partial f^a(\theta_2,\theta_1)}{\partial\theta_2^b}
    h^b{}_{c}\bigl(f(\theta_2,\theta_1)\bigr)
    =h^a{}_{c}(\theta_2).
  \tag{2.B.7}\label{eq:2.B.7}
\end{equation}
Along the second segment, the differential equation~\eqref{eq:2.B.2} for
$U_{\gamma}(s)$ is thus the same as the differential equation for
$U_{\theta_2}(2s-1)$. They satisfy different initial conditions, but
$U_{\gamma}(s)U^{-1}_{\theta_1}(1)$ also satisfies the same differential
equation as $U_{\theta_2}(2s-1)$, and in addition the same initial
condition: at $s=\tfrac12$, both are unity. We therefore conclude that for
$\tfrac12\leq s\leq1$,
\begin{equation*}
  U_{\gamma}(s)U^{-1}_{\theta_1}(1)=U_{\theta_2}(2s-1)
\end{equation*}
and in particular
\begin{equation}
  U_{\gamma}(1)=U_{\theta_2}(1)U_{\theta_1}(1).
  \tag{2.B.8}\label{eq:2.B.8}
\end{equation}
However, this does \emph{not} say that $U_{\theta}(1)$ satisfies the desired
composition rule \eqref{eq:2.B.1}, because although the path
$\Theta_{\gamma}(s)$ runs from $\theta^a=0$ to
$\theta^a=f^a(\theta_2,\theta_1)$, in general it will not be the same as
whatever `standard' path $\Theta_{f(\theta_2,\theta_1)}$ we have chosen to
run directly from $\theta^a=0$ to
$\theta^a=f^a(\theta_2,\theta_1)$. We need to show that $U_{\theta}(1)$ is
independent of the path from 0 to $\theta$ in order to be able to identify
$U[\theta]$ as $U_{\theta}(1)$.

For this purpose, consider the variation $\delta U$ of $U_{\theta}(s)$
produced by a variation $\delta\Theta(s)$ in the path from 0 to $\theta$.
Taking the variation of Eq.~\eqref{eq:2.B.2} gives the differential equation
\begin{align*}
  \frac{d}{ds}\delta U
  &=it_a\delta U\,h^a{}_{b}(\Theta)\frac{d\Theta^b}{ds}
    +it_aU h^a{}_{b,c}(\Theta)\delta\Theta^c
      \frac{d\Theta^b}{ds}\\
  &\quad+it_aU h^a{}_{b}(\Theta)\frac{d\delta\Theta^b}{ds},
\end{align*}
where $h^a{}_{b,c}\equiv\partial h^a{}_{b}/\partial\theta^c$. Using the
Lie commutation relations \eqref{eq:2.2.22} (without central charges) and rearranging
a bit, this gives
\begin{align}
  \frac{d}{ds}\bigl(U^{-1}\delta U\bigr)
  &=\frac{d}{ds}\bigl(iU^{-1}t_aU h^a{}_{b}\delta\Theta^b\bigr)
  \notag\\
  &\quad+iU^{-1}t_aU\delta\Theta^b\frac{d\Theta^c}{ds}
    \left(h^a{}_{c,b}-h^a{}_{b,c}
      +C^a{}_{ed}h^e{}_{b}h^d{}_{c}\right).
  \tag{2.B.9}\label{eq:2.B.9}
\end{align}
However, by taking the limit $\theta_3,\theta_2\to0$ in the associativity
condition \eqref{eq:2.B.6}, we find for all $\theta$:
\begin{equation}
  h(\theta)^a{}_{b,c}
    =-f^a{}_{de}h(\theta)^d{}_{b}h(\theta)^e{}_{c},
  \tag{2.B.10}\label{eq:2.B.10}
\end{equation}
where $f^a{}_{de}$ is the coefficient defined by \eqref{eq:2.2.19}.
Antisymmetrizing in $b$ and $c$ shows that the last term in Eq.~\eqref{eq:2.B.9}
vanishes
\begin{equation}
  h^a{}_{c,b}-h^a{}_{b,c}
    +C^a{}_{ed}h^e{}_{b}h^d{}_{c}=0.
  \tag{2.B.11}\label{eq:2.B.11}
\end{equation}
Eq.~\eqref{eq:2.B.9} thus tells us that the quantity
\begin{equation*}
  U^{-1}\delta U-iU^{-1}t_aU h^a{}_{b}\delta\Theta^b
\end{equation*}
is constant along the path $\theta(s)$. It follows that $U_{\theta}(1)$ is
stationary under any infinitesimal variation of the path that leaves the
endpoints $\Theta(0)=0$ and $\Theta(1)=\theta$ (and
$U_{\theta}(0)=\mathbf{1}$) fixed. But assumption (b) tells us that any path
from $\Theta(0)=0$ to $\Theta(1)=\theta$ can be continuously deformed into
any other, so we may now regard $U_{\theta}(1)$ as a path-independent
function of $\theta$ alone:
\begin{equation}
  U_{\theta}(1)\equiv U[\theta].
  \tag{2.B.12}\label{eq:2.B.12}
\end{equation}
In particular, since the path $\gamma$ leads from 0 to
$f(\theta_2,\theta_1)$, we have
\begin{equation}
  U_{\gamma}(1)=U\bigl[f(\theta_2,\theta_1)\bigr]
  \tag{2.B.13}\label{eq:2.B.13}
\end{equation}
so that Eq.~\eqref{eq:2.B.8} shows that $U[\theta]$ satisfies the group
multiplication law \eqref{eq:2.B.1}, as was to be proved.

Now that we have constructed a non-projective representation $U[\theta]$,
it remains to prove that any projective representation
$\widetilde U[\theta]$ of the same group with the same representation
generators $t_a$ can only differ from $U[\theta]$ by a phase:
\begin{equation*}
  \widetilde U[\theta]=e^{i\alpha(\theta)}U[\theta]
\end{equation*}
so that the phase $\phi$ in the multiplication law for
$\widetilde U[\theta]$:
\begin{equation*}
  \widetilde U[\theta']\widetilde U[\theta]
    =e^{i\phi(\theta',\theta)}
      \widetilde U\bigl[f(\theta',\theta)\bigr]
\end{equation*}
can be removed by a simple change of phase of $\widetilde U[\theta]$. To
see this, consider the operator
\begin{equation*}
  U[\theta]^{-1}U[\theta']^{-1}
    \widetilde U[\theta']\widetilde U[\theta]
  =U\bigl[f(\theta',\theta)\bigr]^{-1}
    \widetilde U\bigl[f(\theta',\theta)\bigr]
    e^{i\phi(\theta',\theta)}.
\end{equation*}
Because $U[\theta]$ and $\widetilde U[\theta]$ have the same generators,
the derivative of the left-hand side with respect to $\theta'{}^a$ vanishes
at $\theta'=0$, and so
\begin{equation*}
  0=\frac{\partial}{\partial\theta^b}
    \left\{U[\theta]^{-1}\widetilde U[\theta]\right\}
    +i\phi_b(\theta)U[\theta]^{-1}\widetilde U[\theta],
\end{equation*}
where
\begin{equation*}
  \phi_b(\theta)\equiv h^a{}_{b}(\theta)
    \left[
      \frac{\partial}{\partial\theta'{}^a}\phi(\theta',\theta)
    \right]_{\theta'=0}.
\end{equation*}
Differentiating this result with respect to $\theta^c$ and
antisymmetrizing in $b$ and $c$ gives immediately
\begin{equation*}
  0=\frac{\partial\phi_b(\theta)}{\partial\theta^c}
    -\frac{\partial\phi_c(\theta)}{\partial\theta^b}.
\end{equation*}
A familiar theorem~\hyperref[ref:2.13]{[13]} tells us that in a simply connected
space, this requires that $\phi_b$ is just a gradient of some function
$\beta$:
\begin{equation*}
  \phi_b(\theta)=\frac{\partial\beta(\theta)}{\partial\theta^b}.
\end{equation*}
Thus the quantity
$U[\theta]^{-1}\widetilde U[\theta]e^{i\beta(\theta)}$ is actually constant
in $\theta$. Setting it equal to its value at $\theta=0$, we see that
$\widetilde U$ is just proportional to $U$:
\begin{equation*}
  \widetilde U[\theta]
    =U[\theta]\exp\bigl(-i\beta(\theta)+i\beta(0)\bigr)
\end{equation*}
as claimed above.

\begin{center}
$*\ *\ *$
\end{center}

The above analysis provides some information about the nature of the phase
factors that can appear in the group multiplication law when the Lie
algebra is free of central charges but the group is not simply connected.
Suppose that the path $\gamma$ from zero to $\theta$ to
$f(\bar\theta,\theta)$ cannot be deformed into the standard path we have
chosen to go from zero to $f(\bar\theta,\theta)$, or in other words, that
the loop from zero to $\theta$ to $f(\bar\theta,\theta)$ and then back to
zero is not continuously deformable to a point. Then
$U[f(\theta_2,\theta_1)]^{-1}U[\theta_2]U[\theta_1]$ can be a phase factor
$\exp(i\phi(\theta_2,\theta_1))\neq\mathbf{1}$, but $\phi$ will be the same
for all other loops into which this can be continuously deformed. The set
consisting of all loops that start and end at the origin and that can be
continuously deformed into a given loop is known as the \emph{homotopy
class}~\hyperref[ref:2.14]{[14]} of that loop; we have thus seen that
$\phi(\theta_2,\theta_1)$ depends only on the homotopy class of the loop
from zero to $\theta$ to $f(\bar\theta,\theta)$ and then back to zero. The
set of homotopy classes forms a group; the `product' of the homotopy class
for loops $\mathcal{L}_1$ and $\mathcal{L}_2$ is the homotopy class of the
loop formed by going around $\mathcal{L}_1$ and then $\mathcal{L}_2$; the
`inverse' of the homotopy class of the loop $\mathcal{L}$ is the homotopy
class of the loop obtained by going around $\mathcal{L}$ in the opposite
direction; and the `identity' is the homotopy class of loops that can be
deformed into a point at the origin. This group is known as the
\emph{first homotopy group} or \emph{fundamental group} of the space in
question. It is easy to see that the phase factors form a representation
of this group: if going around loop $\mathcal{L}$ gives a phase factor
$e^{i\phi}$, and going around loop $\bar{\mathcal{L}}$ gives a phase factor
$e^{i\bar\phi}$, then going around both loops gives a phase factor
$e^{i\phi}e^{i\bar\phi}$. Hence we can catalog all the possible types of
projective representations of a given group $\mathcal{G}$ (with no central
charges) if we know the one-dimensional representations of the first
homotopy group of the parameter space of $\mathcal{G}$. Homotopy groups
will be discussed in greater detail in Volume II.
